\documentclass[12pt]{article}

\usepackage{sbc-template}
\usepackage[brazil]{babel}
\usepackage{graphicx}
\usepackage{times}
\usepackage{url}

\sloppy

\title{Atualização, Ética e Memória na Terapia Cognitivo-Comportamental: Uma Leitura Pessoal do \emph{Procognitiva} (Ciclo 1, Volume 1)}

\author{Yago Melo da Costa}

\begin{document}
\maketitle

\begin{resumo}
Este artigo apresenta uma análise diagnóstica do volume 1 do Ciclo 1 do \emph{Procognitiva: Programa de Atualização em Terapia Cognitivo-Comportamental} (Artmed, 2014), com base no acesso possível ao programa: a amostra oficial de duas páginas --- abertura do capítulo ``Exercícios físicos no contexto da Terapia Cognitivo-Comportamental'', de Aline Sardinha e Antonio Egidio Nardi --- e a ficha editorial. Demonstra-se que o fragmento já é diagnóstico do programa e do campo: a introdução administra publicamente os limites da terapêutica isolada; os objetivos desenham a cartografia do que a TCC hoje reconhece como legítimo; e o muro de acesso inscrito no próprio documento expõe a forma mercadoria da atualização profissional. O artigo conclui com reflexões pessoais e conexões com Kafka, Saramago, Krenak e a \emph{Breve História da Psicologia Moderna Ocidental}.
\end{resumo}

\begin{abstract}
This article presents a diagnostic analysis of volume 1 of Cycle 1 of \emph{Procognitiva: Continuing Update Program in Cognitive-Behavioral Therapy} (Artmed, 2014), based on the access actually available: the official two-page sample --- the opening of the chapter ``Physical exercises in the context of Cognitive-Behavioral Therapy'', by Aline Sardinha and Antonio Egidio Nardi --- and the publisher's catalog entry. It shows that the fragment is already diagnostic of the program and of the field: the introduction publicly administers the limits of isolated therapies; the objectives map what CBT currently recognizes as legitimate; and the paywall embedded in the document itself exposes the commodity form of professional updating. The article closes with personal reflections and connections to Kafka, Saramago, Krenak, and the \emph{Brief History of Modern Western Psychology}.
\end{abstract}

\section{Introdução}

\noindent
\textbf{Acesso e edição utilizada:} O \emph{Procognitiva: Programa de Atualização em Terapia Cognitivo-Comportamental}, Ciclo 1, Volume 1, foi publicado pela Artmed Panamericana em 2014, com 151 páginas, organizado por Carmem Beatriz Neufeld, Eliane Mary de Oliveira Falcone e Bernard Rangé \cite{procognitiva2014}. O volume integral não foi lido para este ensaio: o acesso obtido corresponde à amostra oficial de duas páginas, contendo a introdução e os objetivos do capítulo ``Exercícios físicos no contexto da Terapia Cognitivo-Comportamental'', de Aline Sardinha e Antonio Egidio Nardi \cite{sardinha-nardi}, além da ficha editorial. Como a amostra é material de web, sem foliação interna, as citações são textuais, sem números de página. A restrição não é preâmbulo defensivo: é matéria do artigo --- o que um fragmento autoriza a dizer, e o que não autoriza.

O \emph{Procognitiva} existe para uma profissão que precisa se atualizar sem parar, e o volume 1 do Ciclo 1 é o primeiro degrau desse programa de educação continuada à distância. Este artigo pergunta o que o programa \emph{ensina} para além de seus temas --- e a resposta, já visível no fragmento, é tripla. O campo se define hoje pelas interfaces entre saberes, não por fronteiras (seção~3); se mede por verbos de cartografia, não de conversão (seção~4); e se transmite sob forma mercadoria, com o muro de acesso parte integrante do próprio documento (seção~4). A seção~2 examina o que a escolha do capítulo-vitrine revela do programa; a seção~5 reúne as reflexões pessoais e as conexões com as demais leituras da série.

\section{O programa e o que o fragmento revela}

O volume não é livro de autor único: é coletânea organizada, e a organização é o primeiro argumento. Neufeld, Falcone e Rangé \cite{procognitiva2014} não assinam capítulos --- assinam a costura entre eles, garantindo que especialistas distintos falem a mesma língua didática. O \emph{Procognitiva} é, nesse sentido, menos uma obra do que um \emph{sistema}: ciclos anuais, volumes temáticos, e um leitor profissional esperado ao fim de cada capítulo. A educação continuada à distância, aqui, não é meio: é formato de existência profissional --- atualizar-se é o que o psicólogo \emph{é}, não o que faz ocasionalmente. Há, nisso, uma teoria implícita do tempo profissional: o ano é a unidade da obsolescência, e o volume é a unidade da cura. Os 151 páginas não são leitura --- são dosagem: o suficiente para reabastecer um ano de prática até o próximo ciclo abrir. E a coletânea administra um duplo vínculo elegante: cada capítulo carrega a autoridade de especialistas nomeados (Sardinha, Nardi e os demais), enquanto o programa inteiro permanece sem rosto único --- a marca é o sistema, não o autor. Quem assina é quem vende; quem escreve é quem credencia.

A amostra confirma o diagnóstico pelo que escolhe mostrar. O capítulo-vitrine é ``Exercícios físicos no contexto da Terapia Cognitivo-Comportamental'', de Sardinha e Nardi \cite{sardinha-nardi} --- o corpo entrando pela porta da terapia cognitiva. A escolha não é arbitrária: num programa que vende atualização, a vitrine precisa exibir o que há de mais currency no campo, e a integração entre saúde física e mental é, no próprio texto, o vetor declarado. O capítulo promete ensinar ao clínico a enxergar a interface entre o movimento do corpo e o sofrimento psíquico --- entre o sedentário e o ansioso, o deprimido e o parado.

O que o fragmento revela, antes de qualquer conteúdo, é o \emph{contrato didático} do programa. Mesmo truncado pela paga, o documento anuncia sua estrutura esquelética: uma Introdução que justifica o tema; uma lista de objetivos introduzida pela fórmula ``Ao final do artigo, o leitor poderá:''; e, em seguida, o cabeçalho ``Esquema conceitual'', prometendo um mapa antes do argumento. Objetivos antes de conteúdo, mapa antes de discussão, referências ao fim: a máquina didática se apresenta antes de funcionar. O leitor é treinado não apenas em TCC, mas num \emph{gênero} --- o capítulo de atualização, com seus tempos e suas promessas. E o gênero, ele próprio, é doutrina: a ideia de que conhecimento clínico chega em porções numeradas, consumíveis entre plantões.

\section{A introdução como manifesto da interdisciplinaridade}

A primeira página do capítulo abre com uma declaração de princípios que merece ser lida com o cuidado que se dá a um manifesto: ``Nos últimos anos, tem se tornado cada vez mais patente o entendimento de que a integração de conhecimentos e o trabalho multidisciplinar, interdisciplinar e transdisciplinar devem embasar as ações de prevenção e tratamento em saúde''. A escalada dos prefixos --- multi, inter, trans --- não é preciosismo: é a ambição declarada em gramática. Cada prefixo promete um grau a mais de dissolução das fronteiras entre saberes, e a frase os empilha como quem empilha credenciais.

O parágrafo seguinte é o mais corajoso do fragmento, e talvez do gênero. Em vez de anunciar triunfos, o texto administra perdas: ``Em razão das limitadas taxas de eficácia alcançadas pelas diversas modalidades de tratamento isoladamente (psicofármacos, psicoterapias, terapias alternativas etc.), é crescente a conscientização de que o padrão-ouro de terapêutica dos diversos transtornos mentais perpassa várias esferas do conhecimento''. Três movimentos nessa frase. Primeiro, a admissão pública de que nada, sozinho, funciona bem: ``limitadas taxas de eficácia''. Segundo, o parêntese que inventaria --- psicofármacos, psicoterapias, terapias alternativas --- nivelando farmácura e terapia na mesma fila de insuficiências, sem hierarquia prévia. Terceiro, o destino do adjetivo ``padrão-ouro'': ele não recai mais sobre uma técnica, mas sobre uma \emph{combinação} --- ``perpassa várias esferas do conhecimento''. O ouro mudou de substância para alquimia: não está no método que se aplica, está no arranjo que se orquestra.

Dessa base, o capítulo declara seu objeto com precisão de instrução de uso: ``Este artigo visa a sistematizar as situações mais comuns em saúde mental, nas quais a integração dos conhecimentos da psicoterapia cognitivo-comportamental com os profissionais prescritores de exercícios físicos se faz necessária''. A clínica aparece como negociação entre prescrições: o psicólogo do texto não prescreve exercício --- encaminha, para ``utilizar o encaminhamento para exercício físico supervisionado como ferramenta em sua prática clínica''. O instrumento é burocrático por um lado (um encaminhamento) e corporal por outro (exercício supervisionado); a ferramenta da alma moderna, mais uma vez, é um formulário que põe o corpo em movimento. O parágrafo fecha com a aritmética da sinergia: ``a combinação de terapias potencialize os resultados do tratamento e a qualidade de vida dos pacientes de maneira mais eficaz do que a soma das intervenções isoladas'' --- 1+1 maior que 2, elevado a princípio terapêutico.

O diagnóstico que se tira da introdução é o de um campo que se define pelas \emph{interfaces}. A psicologia representada neste capítulo não guarda fronteiras, não reivindica impérios de objeto: ela se apresenta como nó num tecido de saberes, e sua competência é articular --- médicos, psicólogos, nutricionistas, educadores físicos, fisioterapeutas, terapeutas ocupacionais, odontólogos, fonoaudiólogos aparecem na lista de interlocutores como plural obrigatório do cuidado. É uma psicologia que, em sua melhor formulação pública, começa confessando o que não sabe fazer sozinha. Note-se ainda o vocabulário do quadro: o campo da ação é chamado, sem cerimônia, de ``saúde mental'' --- a casa secular da psicologia é a medicina, e é dentro dela que o capítulo negocia sua autoridade. E o verbo do manifesto é ``deve'' --- ``devem embasar as ações de prevenção e tratamento'' ---, imperativo de profissão que legisla sobre si mesma: a interdisciplinaridade, aqui, não é convite aberto, é standard declarado, e o capítulo inteiro se apresenta como sua demonstração clínica.

\section{Os objetivos e a anatomia do contrato didático}

Os objetivos do capítulo merecem leitura palavra a palavra, porque são o retrato do que o campo se dispõe a ensinar. ``Ao final do artigo, o leitor poderá'': ver ``como a ansiedade pode afetar a prática de exercícios físicos, contribuindo para um estilo de vida sedentário''; ``identificar o papel dos exercícios físicos no tratamento da ansiedade e do transtorno de pânico''; ``estabelecer a relação entre sedentarismo, ansiedade e saúde física''; ``identificar o papel dos exercícios físicos no tratamento da depressão''; ``reconhecer comportamentos disfuncionais dos pacientes relacionados à prática de exercícios físicos que podem ser indicativos de aspectos psicopatológicos''; ``identificar como o exercício pode ser utilizado de forma compulsiva por pacientes com alterações da autoimagem corporal''; e ``estabelecer um plano de tratamento conjunto com profissionais prescritores de exercício'' \cite{sardinha-nardi}.

Duas observações de anatomia. A primeira é o vocabulário dos verbos: ver, identificar (três vezes), estabelecer (duas), reconhecer. Nenhum objetivo diz ``convencer'', ``converter'' ou ``salvar'': o contrato é cartográfico --- treinar a percepção do clínico, não instalar doutrina. A psicologia deste capítulo quer olhos, não prosélitos. A segunda observação é o conteúdo do mapa: ele inclui a face escura da própria ferramenta. Dois objetivos dedicam-se a mapear o exercício como vetor de patologia --- compulsão, autoimagem corporal distorcida, comportamentos que ``podem ser indicativos de aspectos psicopatológicos''. Um campo que descreve os riscos do seu instrumento no mesmo capítulo em que o recomenda pratica uma espécie de vacinação contra o próprio entusiasmo. E o último objetivo não é clínico, mas institucional: ``estabelecer um plano de tratamento conjunto'' --- a interface de novo, agora como entregável: a cooperativa entre saberes ensinada como tarefa, não como ideal.

A ordem do mapa também é diagnóstica. O capítulo abre pela ansiedade --- com o transtorno de pânico nomeado ao lado ---, e logo depois vem a depressão: as duas grandes narrativas da saúde mental contemporânea recebem o corpo como porta de entrada. Entre elas aparece o ``estilo de vida sedentário'', e a expressão merece atenção: um hábito social, nascido de cidades, jornadas e desigualdades, é traduzido para o registro clínico como variável a ser intervida no indivíduo. O social não desaparece --- é encapsulado no paciente que se recusa a andar. Sete objetivos, nenhum deles prometendo cura: a humildade da cartografia é coerente com a humildade da introdução. O que se promete é ver melhor; o que se vende, no fim das contas, é atenção treinada.

Depois dos objetivos, o documento anuncia ``Esquema conceitual'' --- e termina. As palavras finais da amostra não são do autor: são do muro. ``Cadastre-se para ler o artigo completo''. O corte é editorial e econômico ao mesmo tempo, e o diagnóstico que dele se tira é o mais incômodo deste artigo: no \emph{Procognitiva}, a atualização é mercadoria. O programa que existe para impedir que a psicologia repita seus erros de amnésia existe, ele próprio, sob a forma do pagamento --- ciclos, volumes, assinaturas. Não há ironia acidental nisso: há a forma social do conhecimento profissional em 2014, exibida sem disfarce. O que se vende não é informação --- está em toda parte ---, é \emph{atualização credenciada}: o carimbo de quem estudou este ano. E a amostra é honesta o bastante para exibir a porta em vez de fingir que não existe: ela não devolve o conteúdo de graça e pede desculpas --- ela informa, com clareza comercial, que o acesso tem preço. É o resto do documento que se recusa a ser lido, e essa recusa, inscrita no próprio arquivo, é a página mais sociológica da amostra.

\section{Reflexões pessoais e conexões com as demais leituras}

Li duas páginas de um livro de cento e quinze mais quinze. E foi suficiente para aprender o que a leitura de fragmentos, quando honesta, pode ensinar: o retrato de um campo no momento em que ele se olha no espelho --- e a diferença entre o espelho e o resto da casa. A promessa do capítulo, tal como o fragmento a deixa entrever, é a de uma psicologia que começa por dizer o que não faz sozinha. Se o restante dos 151 páginas mantém esse registro, o programa cumpre o que anuncia; se não mantém, o muro teve o duplo efeito de vender e de proteger. Não sei qual dos dois livros existe atrás da porta --- e essa não-saber é a posição ética que a amostra me impôs, e que prefiro a simulação de uma leitura que não fiz.

Com Kafka \cite{kafka2022}, o muro tem parentesco burocrático: ``Cadastre-se para ler o artigo completo'' é a porta da recepção de Gregor, o formulário que decide quem entra e quem é processado --- e o conhecimento profissional, aqui, tem o mesmo fluxo que o do conto: quem não preenche, não aparece. A diferença importa, porém: a burocracia de Kafka oculta que é burocracia; a do \emph{Procognitiva} exibe-se como porta. Há mais decência numa porta anunciada do que numa porta que finge ser paisagem. Com Saramago \cite{saramago1997}, o caso é o do mapa sem a ilha: a ficha editorial promete os ciclos, os volumes, os temas --- todas as ilhas estão nos mapas do catálogo ---, e a ilha propriamente dita, o capítulo, permanece do outro lado da água. Os marinheiros do rei ficaram com as ilhas conhecidas; eu fiquei com as duas páginas conhecidas, e a caravela não partiu por falta de moeda, não por falta de mar.

Com Krenak \cite{krenak2019}, a conexão é o corpo. O capítulo-vitrine dedica-se a devolver movimento a gente ``totalmente alienada do mínimo exercício de ser'' --- a fórmula do diagnóstico krenakiano encontra aqui seu contraponto clínico: onde Krenak denuncia a alienação do ser, Sardinha e Nardi desenham o encaminhamento que devolve o exercício ao sedentário. São duas descrições do mesmo paciente, uma escrita de dentro de uma aldeia, outra de dentro de um ambulatório. Lidas juntas, as duas páginas da amostra ganham uma terceira dimensão: o sedentário do capítulo é, também, o produto do ``liquidificador chamado humanidade'' que Krenak diagnostica \cite[p. 14]{krenak2019} --- o corpo arrancado de seus coletivos e reencaixado, depois, por encaminhamento médico, numa esteira que devolve o movimento sem devolver o pertencimento. Se a TCC quiser levar a sério sua própria admisão de ``limitadas taxas de eficácia'', talvez precise ler Krenak como consulta externa: exercício devolve o corpo ao ser; história devolve o ser ao mundo.

E com a \emph{Breve História da Psicologia Moderna Ocidental}, de Teresa Sousa Machado \cite{machado}, o \emph{Procognitiva} forma o par que a série inteira pedia: o livro de história mostra o que acontece quando a psicologia esquece suas escolhas --- eugenia, testes ao serviço da exclusão, experimentos sem ética ---; o programa de atualização é, ele mesmo, a institucionalização do não-esquecer: ciclos anuais contra a amnésia profissional. A frase que Machado repete de Ebbinghaus --- a psicologia tem ``um longo passado e uma curta história'' \cite[p. 1]{machado} --- encontra no \emph{Procognitiva} sua tradução administrativa: para que a história curta não encolha ainda mais, publica-se um capítulo por mês. Resta a pergunta que o muro devolve: a memória de uma profissão pode ser institucionalizada como assinatura? A amostra diz que, em 2014, era exatamente assim que se fazia.

O fragmento me deixou, por fim, duas disciplinas. A primeira é a da honestidade de escala: declarar de quantas páginas se fala, e não fingir o volume inteiro na voz do resenha --- o que se sabe, o que se supõe e o que se paga devem constar do mesmo parágrafo. A segunda é a da escuta do espelho: o que este campo diz de si em seu melhor momento --- ``limitadas taxas de eficácia'' --- vale mais, como diagnóstico, do que o que ele diz de si em seus compêndios de glória, porque as frases de glória qualquer instituição sabe fabricar, e as de perda só se escrevem sob pressão da verdade. Se a atualização é uma forma de memória ativa, então ler duas páginas com rigor já é participar do ritual: contar mais uma história, ainda que a história tenha o tamanho exato do que se pode pagar --- e perguntar, do outro lado da porta, o que custa a porta aberta.

\begin{thebibliography}{99}

\bibitem{procognitiva2014}
NEUFELD, Carmem Beatriz; FALCONE, Eliane Mary de Oliveira; RANGÉ, Bernard (Orgs.).
\emph{Procognitiva: Programa de Atualização em Terapia Cognitivo-Comportamental. Ciclo 1, Volume 1}.
Porto Alegre: Artmed Panamericana, 2014. 151 p.

\bibitem{sardinha-nardi}
SARDINHA, Aline; NARDI, Antonio Egidio.
Exercícios físicos no contexto da Terapia Cognitivo-Comportamental.
In: \emph{Procognitiva: Programa de Atualização em Terapia Cognitivo-Comportamental. Ciclo 1, Volume 1}.
Porto Alegre: Artmed Panamericana, 2014.
Amostra oficial digitalizada (2 p.), publicada pela editora.

\bibitem{kafka2022}
KAFKA, Franz.
\emph{A Metamorfose} [\emph{Die Verwandlung}, 1915].
Tradução de Giovane Rodrigues.
São Paulo: Instituto Mojo, 2022. (Coleção Literatura Livre).
Edição bilíngue Português-Alemão. E-book.
ISBN 978-65-89008-23-1.
Disponível em: \url{https://mojo.org.br/wp-content/uploads/2024/08/006-A-metamorfose-978-65-89008-23-1.pdf}.

\bibitem{saramago1997}
SARAMAGO, José.
\emph{O Conto da Ilha Desconhecida}.
São Paulo: Companhia das Letras, 1997.

\bibitem{krenak2019}
KRENAK, Ailton.
\emph{Ideias para adiar o fim do mundo}.
1ª ed. São Paulo: Companhia das Letras, 2019.
ISBN 978-85-359-3241-6.

\bibitem{machado}
MACHADO, Teresa Sousa.
\emph{Breve História da Psicologia Moderna Ocidental}.
Documento PDF, 128 p.

\end{thebibliography}

\end{document}
